\documentclass[conference]{IEEEtran}
\usepackage{iftex}
\ifPDFTeX\usepackage[utf8]{inputenc}\usepackage[T1]{fontenc}\fi
\usepackage{amsmath,amssymb,bm,mathtools,newtxtext,newtxmath}
\usepackage{booktabs,array,graphicx,adjustbox,placeins,siunitx}
\usepackage{pgfplots}\pgfplotsset{compat=1.18}
\usepackage[final]{microtype}
\usepackage{cite}
\usepackage[breaklinks=true,hidelinks]{hyperref}
\usepackage[nameinlink,capitalise,noabbrev]{cleveref}
\interdisplaylinepenalty=2500\emergencystretch=2em\allowdisplaybreaks\raggedbottom
\setcounter{topnumber}{5}\setcounter{dbltopnumber}{5}
\renewcommand{\topfraction}{.95}\renewcommand{\dbltopfraction}{.95}
\renewcommand{\textfraction}{.05}\renewcommand{\floatpagefraction}{.8}
\renewcommand{\dblfloatpagefraction}{.8}
\makeatletter
\setlength{\@fptop}{0pt}\setlength{\@dblfptop}{0pt}
\setlength{\@fpsep}{12pt plus 2pt}\setlength{\@dblfpsep}{12pt plus 2pt}
\makeatother
\newcommand{\R}{\mathbb R}
\newcommand{\vc}[1]{\bm{#1}}
\newcommand{\sk}[1]{[#1]_\times}
\newcommand{\tr}{^{\mathsf T}}
\newcommand{\IncludePGF}[1]{\resizebox{\linewidth}{!}{\input{#1}}}
\providecommand{\mathdefault}[1]{#1}
\hypersetup{pdftitle={Tangent-Bias Geometry and Bias Learning in Shared-Mode CARMA Marine Estimation},pdfauthor={Mikhail S. Grushinskiy}}
\input{tangent-results-generated.tex-part}
\title{Tangent-Bias Geometry and Bias Learning\\in Shared-Mode CARMA Marine Estimation}
\author{\IEEEauthorblockN{Mikhail S. Grushinskiy}
\IEEEauthorblockA{Bareboat Necessities, ocean-IMU project\\Independent Researcher, Fair Lawn, NJ, USA}}
\begin{document}\maketitle
\begin{abstract}
\normalfont
Tangent-group inertial filters geometrically couple sensor biases with
navigation states, but their advantages depend on the underlying dynamics
and available aiding. We investigate a vessel-mounted accelerometer,
gyroscope and full three-axis magnetometer without external velocity or
position measurements. The existing minimal shared-mode CARMA(4,1) prior
models wave-relative motion. We implement a 27-state tangent-bias hybrid on
$T\mathrm{SE}(3)$ extended by Euclidean wave coordinates, preserving the
current process model, startup, measured-endpoint gyro integration,
adaptation and covariance writes. The full tangent group of
$\mathrm{SE}_2(3)$ instead introduces an additional virtual bias; its exact
standard-INS navigation-error result is not asserted for this CARMA hybrid.
A shared directional bias-gain veto tests a conditional excitation monitor
with full Joseph covariance updates. Matched sensor histories separate
geometry, bias weighting, startup and live recovery. \StudyHeadline
Nonlinear sensor symmetries and nuisance-eliminated window information
explain why magnetic aiding and fixed-heading roll/pitch excitation can
help, while low-frequency motion and reference uncertainty retain weak
bias directions. Results are experimental and refer to wave coordinates,
not absolute geographic navigation.
\end{abstract}
\begin{IEEEkeywords}
Equivariant filtering, tangent group, inertial sensor bias, CARMA, marine
motion, observability, statistical consistency.
\end{IEEEkeywords}

\section{Question, Baseline and Research Scope}
The practical question is whether a different geometric error and a more
cautious learning policy improve live motion and bias estimates. Mathematical
elegance alone is insufficient. The baseline is main commit
\texttt{7712a0e046bb653839a4d582ec893d7671ac1e5f}, including the November 2026
measured-endpoint gyro integration, three-dimensional bias corrections,
covariance fixes and QMEKF sample-path optimization. Comparisons against
historical endpoint prediction would confound this experiment.
Integration also preserves main revision \texttt{480d78f}, which tightens
production replay limits and rejects invalid scored histories. A fresh
ten-filter, 900 s wave run after that integration reproduces every cohort
score and sensor hash exactly, excluding host runtime and object size.

The sensors observe specific force, angular rate and magnetic field. Outputs
$\vc p,\vc v,\vc a$ describe oscillatory displacement about a local wave
equilibrium. Neither absolute position, constant translation nor current is
identified by these sensors. Position-drift metrics below are explicitly
wave-coordinate errors. The shipping OU-III estimator is unchanged.

\subsection{Verified literature and its transfer limits}
The MEKF treats attitude multiplicatively and other states additively.
An invariant EKF on $\mathrm{SE}_2(3)$ geometrically represents attitude,
velocity and position; adding Euclidean IMU biases generally loses the
bias-free group-affine property. The two-frame group (TFG) incorporates
world vectors and body-frame biases in one asymmetric group law
\cite{barrau2023}. Fornasier et al. introduced the tangent-group bias
extension, including a virtual input and virtual bias, for standard biased
inertial navigation \cite{fornasier2022}. The 2025 comparison interprets
MEKF, imperfect IEKF, TFG, full TG, direct-position and semi-direct-bias
filters as symmetry choices \cite{fornasier2025}. Its favorable full-TG
result concerns externally aided trajectories with sufficient observability;
it is not evidence about an IMU/magnetometer wave estimator.

The 2024 dissertation gives a broader geometric treatment
\cite{fornasier2024}. The September 2026 IterEqF preprint treats an iterated
correction as Gauss--Newton on a manifold and derives its covariance reset
\cite{vantoor2026}. Its reported localization example is not this marine
sensor suite. A separate September 2026 acoustic/depth preprint studies
tangent-group filtering with DVL velocity and pressure aiding
\cite{lunawat2026}. These additional measurements carry independent motion
information absent here. Both 2026 works are cited as preprints; their
published abstracts motivate future comparisons rather than establish an
accuracy claim for ocean-IMU.

\subsection{Which tangent group is implemented?}
The standard full TG-INS uses
$\mathrm{SE}_2(3)\ltimes\mathfrak{se}_2(3)$: nine navigation coordinates
and nine bias coordinates, including a virtual position-input bias. Adding
the twelve remaining minimal CARMA coordinates gives 30 errors. Fixing that
virtual bias to zero without changing the error/reset equations does not
preserve the full symmetry.

We instead use the six-bias semi-direct construction whose velocity/bias
subgroup is $T\mathrm{SE}(3)$, with the other wave vectors as Euclidean
representations. This preserves 27 states. We call it \emph{tangent-bias
hybrid}, abbreviated TAN in tables. It is not a full 30-state TG-EqF
implementation or a demonstration of its autonomous standard-INS error
theorem. This distinction is a deliberate model-preserving design choice,
and limits what a negative comparison can conclude about full TG filters.

\section{Physical and Stochastic Model}
The body-to-world attitude is $R\in\mathrm{SO}(3)$, with a local NED world
frame and $\vc g=(0,0,g)\tr$. Sensor models are
\begin{align}
 \vc\omega_m&=\vc\omega+\vc b_g+\vc n_g,\\
 \vc f_m&=R\tr(\vc a-\vc g)+\vc b_a+\vc n_a,\label{eq:force}\\
 \vc m_m&=R\tr\vc B+\vc n_m.\label{eq:mag}
\end{align}
Temperature calibration, lever-arm processing and magnetic startup remain
the existing components. Fixed calibration errors and uncertain reference
fields are distinct from white packet noise.

Per spatial axis the existing minimal shared-mode realization has
\begin{equation}
 \begin{aligned}
 \vc x_i&=(p_i,v_i,a_i,\ell_i,w_{i1},w_{i2})\tr,\\
 d\vc x_i&=A(\kappa)\vc x_i\,dt+G(\kappa,s_i)d\vc W_i.
 \end{aligned}
\end{equation}
The first four coordinates are the established CARMA(4,1) realization; the
last two form its lossless coloring factor. The complexified eight-real-state
description contains two unobservable directions per axis. Their removal
preserves the physical spectra, including cross-axis quadrature, and reduces
$3+6+24=33$ to $3+6+18=27$ errors. This established reduction is reused,
not recomputed or approximated by the new geometry.

Let $W=[\vc p,\vc v,\vc a,\vc\ell,\vc w_1,\vc w_2]\in\R^{3\times6}$.
Its nominal dynamics are $\dot W=WA\tr$. Spatial excitation is Hermitian:
for the exact common discretization,
\begin{equation}
 Q_{ij}=\Re(\Sigma_{ij})Q_R+\Im(\Sigma_{ij})Q_I.
\end{equation}
$Q_R$ is symmetric and $Q_I$ antisymmetric; the complete covariance is
real symmetric positive semidefinite. Bias dynamics are the existing random
walk for $\vc b_g$ and active OU prior
$\dot{\vc b}_a=-\lambda\vc b_a+\vc n_{ba}$; a held bias uses $\lambda=0$
and the existing holding policy. No $S=0$ pseudo-measurement is added:
stable CARMA displacement dynamics already regularize low-frequency motion.
That regularization is a model prior, not an external position observation.

\subsection{Spectral fitting and adaptation}
The fixed CARMA spectral fit, period estimator, measurement-only leveling
proxy, amplitude evidence and shared cross-spectrum learner are unchanged.
One common time scale $\kappa$ keeps dynamics isotropic across spatial axes;
amplitudes and cross-axis excitation may differ. Model commits preserve
physical $(p,v,a,\ell)$ and transform normalized coloring states
$\vc w_i/s_i$. The full covariance follows the same physical-coordinate
congruence, retaining navigation/bias/motion cross-covariance. This includes
the attitude shear when different axis transformations fail to commute with
rotation. Reusing the existing commit path avoids confusing an adaptation
coordinate error with geometric filtering performance.

\section{Tangent-Bias Geometry}
\subsection{Group action and exact finite correction}
Let $N=\mathrm{SO}(3)\ltimes\R^{3\times6}$ with
$(R,W)(U,Z)=(RU,W+RZ)$. Projection $\pi(R,W)=(R,\vc v)$ is a homomorphism
to $\mathrm{SE}(3)$. The tangent-bias group is
\begin{equation}
 \mathcal G=N\ltimes_{\operatorname{Ad}_{\pi}}\R^6.\label{eq:group}
\end{equation}
Its bias ideal is an additive vector space carrying the adjoint action;
it is not a second independent rigid-body pose. For
$\vc b=(\vc b_g,\vc b_a)$ define lifted bias
\begin{equation}
 \gamma=\operatorname{Ad}_{(R,\vc v)}\vc b
 =\begin{pmatrix}R\vc b_g\\R\vc b_a+\sk{\vc v}R\vc b_g\end{pmatrix}.
\end{equation}
In physical storage $(R,W,\vc b)$ the group product satisfies
\begin{align}
 R_+&=RU,&W_+&=W+RZ,\\
 \vc b_{g,+}&=U\tr\vc b_g+\vc c_g,\\
 \vc b_{a,+}&=U\tr(\vc b_a-\sk{\vc z_v}\vc b_g)+\vc c_a.
\end{align}
This is implemented as a small extension of the existing two-frame storage.
No CARMA process is copied.

Use a right-invariant error and a left correction:
\begin{equation}
 \eta=\hat X X^{-1},\qquad \hat X^+=\exp_{\mathcal G}(-\hat\xi)\hat X^-.
\end{equation}
To first order, relative to the current two-frame error $\xi_F$,
\begin{align}
 \vc\phi&=\log(\hat R R\tr),\quad
 \vc\rho_i=\hat{\vc w}_i-\hat R R\tr\vc w_i,\\
 \vc\beta_g&\simeq\hat R(\hat{\vc b}_g-\vc b_g),\quad
 \vc\beta_a\simeq\hat R(\hat{\vc b}_a-\vc b_a),\\
 \vc\zeta_a&=\vc\beta_a+\sk{\hat{\vc v}}\vc\beta_g,
 \qquad \xi_T=S(\hat{\vc v})\xi_F.\label{eq:shear}
\end{align}
Only the accelerometer-bias row changes in $S$. For
$u=(\vc\phi,\vc\rho_i,\vc\zeta_g,\vc\zeta_a)$, the exponential uses
\begin{align}
 E&=\exp(\sk{\vc\phi}),\quad W_E=J_l(\vc\phi)\rho,\\
 \gamma_{g,E}&=J_l(\vc\phi)\vc\zeta_g,\\
 \gamma_{a,E}&=J_l(\vc\phi)\vc\zeta_a+
 Q_l(\vc\phi,\vc\rho_v)\vc\zeta_g,
\end{align}
then maps $\gamma_E$ back with
$\operatorname{Ad}_{(E,\vc v_E)}^{-1}$. Here $Q_l$ is the standard
$\mathrm{SE}(3)$ Jacobian cross block, shared with the TFG implementation.
All coloring coordinates retain their original linear stochastic model.

\subsection{Continuous error dynamics and the missing cancellation}
For the repository's two-frame convention let
$\Omega=\sk{\hat R(\vc\omega_m-\hat{\vc b}_g)}$. Ignoring displayed noise
terms, the current linearized process is
\begin{align}
 \dot{\vc\phi}&=-\vc\beta_g,&
 \dot{\vc\beta}_g&=\Omega\vc\beta_g,\\
 \dot{\vc\rho}_i&=\sum_j A_{ij}\vc\rho_j-
                \sk{\hat{\vc w}_i}\vc\beta_g,\label{eq:carmaerror}\\
 \dot{\vc\beta}_a&=(\Omega-\lambda I)\vc\beta_a.
\end{align}
Differentiating \cref{eq:shear} gives
\begin{align}
 F_T&=\dot S S^{-1}+S F_F S^{-1},\label{eq:ft}\\
 \dot{\vc\zeta}_a&=(\Omega-\lambda I)\vc\zeta_a+K_v\vc\beta_g,\\
 K_v&=\sk{\dot{\hat{\vc v}}}+\sk{\hat{\vc v}}\Omega
       -(\Omega-\lambda I)\sk{\hat{\vc v}}.
\end{align}
This retains every first-order coupling. In standard INS, force drives
velocity and the full tangent bias can move the velocity/position
linearization error into bias dynamics. Here force is \emph{measured
against} a latent CARMA acceleration, while \cref{eq:carmaerror} propagates
velocity. Inserting the standard autonomous TG-INS matrix would therefore
change the process model and discard real terms. We do not do so.

\subsection{Discrete propagation and covariance}
The continuous noise injection is also coordinate dependent. Let
$d\vc u_g$ denote gyro measurement noise, $d\vc u_{bg},d\vc u_{ba}$
bias driving noise, and $d\vc u_{w_i}$ the jointly correlated CARMA
driving increments. In the same error convention,
\begin{align}
 d\vc\phi\big|_n&=\hat R\,d\vc u_g,\\
 d\vc\beta_g\big|_n&=-\hat R\,d\vc u_{bg},\\
 d\vc\rho_i\big|_n&=\sk{\hat{\vc w}_i}\hat R\,d\vc u_g-d\vc u_{w_i},\\
 d\vc\beta_a\big|_n&=-\hat R\,d\vc u_{ba}.
\end{align}
These blocks define $L_F$ and
$Q_F^c=L_F Q_{\mathrm{drive}}L_F\tr$, including shared-mode forcing
and gyro-induced attitude/motion cross-covariances. Accelerometer packet
noise belongs to the force likelihood, not a second velocity-driving
channel. Bias driving noise follows the native holding policy.
For a frozen linearization the reference discretization is
\begin{equation}
 \Phi_F=e^{F_Fh},\qquad
 Q_F=\int_0^h e^{F_Fs}Q_F^c e^{F_F\tr s}\,ds.
\end{equation}
The CARMA block evaluates its exact matrix-exponential discretization.
Navigation retains the existing structured measured-endpoint
approximation and shared-endpoint noise bound; it is not replaced by an
assumed autonomous full-TG transition.
The existing exact CARMA discretization and measured-endpoint gyro map
produce $\Phi_F,Q_F$. At corresponding nominal endpoints,
\begin{equation}
 \Phi_T=S_1\Phi_F S_0^{-1},\quad Q_T=S_1Q_F S_1\tr,
 \quad P_T=S P_F S\tr.\label{eq:discrete}
\end{equation}
Thus cross-covariances are retained rather than zeroed. These equations
also show that a first-order coordinate change alone cannot improve the
Gaussian linear update. A measured improvement must come from finite
retraction/reset effects, different higher-order output approximations or
an actual model change.

For efficiency the implementation stores the public covariance in the
current two-frame chart, reusing its optimized prediction and physical
covariance writes. At each correction it maps the posterior to the tangent
chart, applies the tangent reset and maps back. This is algebraically
\cref{eq:discrete}, not an additional filter or an information fusion step.
The existing conservative reused-endpoint gyro-noise treatment is retained;
averaging two samples is not counted as two independent measurements.

\subsection{Measurement model and reset}
In world residual coordinates, without a lever arm,
\begin{align}
 \vc r_a&=\hat R(\vc f_m-\hat{\vc b}_a)-(\hat{\vc a}-\vc g),\\
 H_{a,F}&=[\sk{\vc g}\quad0\quad\cdots\ -I_{\rho_a}\ \cdots\ -I_{\beta_a}],\\
 \vc r_m&=\hat R\vc m_m-\vc B,\quad
 H_{m,F}=[-\sk{\vc B}\quad0\ \cdots].
\end{align}
Noise rotates as $R\mathcal R R\tr$. The tangent Jacobian is
$H_T=H_F S^{-1}$; in particular
$H_{a,T,\beta_g}=\sk{\hat{\vc v}}$ in addition to any existing lever-arm
term. Using this identity through the backing chart yields the same
first-order residual likelihood.

With $d=S_0Kr$, the complete stored-chart reset is
\begin{equation}
 J_{\mathrm{store}}=S_+^{-1}J_{r,\mathcal G}(d)S_0,
 \qquad P^+=J_{\mathrm{store}}P_{\mathrm{post}}J_{\mathrm{store}}\tr.
\end{equation}
The Lie-algebra adjoint gains the blocks
$(\operatorname{ad}_d)_{a,g}=\sk{\rho_v}$ and
$(\operatorname{ad}_d)_{a,v}=\sk{\zeta_g}$. Their two paths also contribute
to the bias/attitude reset block. Omitting either produces an inconsistent
finite reset. A 32-term series is evaluated only on $3\times3$ blocks;
covariance congruences use rank-row operations, not a dense exponential or
an additional $27\times27$ scratch matrix. Finite differences and an
independent dense covariance congruence test verify the implementation.

When a physical bias is held during startup, the original two-frame
retraction is used. Otherwise a second-order tangent cross term could move
a mean whose gain row was held at zero. Switching geometry when native
bias learning becomes active changes no physical state or backing
covariance. This preserves the existing startup semantics and isolates
live geometric effects.

\subsection{Iterated updates: what would be a separate experiment?}
An iterated update would minimize one prior-plus-measurement objective,
\begin{equation}
 \tfrac12\|\epsilon(X,\hat X^-)\|_{P^-}^2+
 \tfrac12\|y-h(X)\|_{\mathcal R}^2,
\end{equation}
relinearizing at trial means and resetting covariance to the final chart.
Reapplying a Kalman update to the same packet as if it were new data would
double-count information. Neither an iterated update nor the full 30-state
TG-INS is silently folded into TAN; their effect cannot be inferred from
the comparisons here.

\section{Bias Observability and Conservative Learning}
\subsection{An exact sensor ambiguity and its limits}
Full three-axis magnetic measurements constrain two attitude directions,
but a rotation $Q$ about a known field obeys $Q\vc B=\vc B$. For any body
bias perturbation $\vc d(t)$ define
\begin{align}
 R'&=QR,\quad \vc b_a'=\vc b_a+\vc d,\\
 \vc a'&=\vc g+Q(\vc a-\vc g-R\vc d).\label{eq:gauge}
\end{align}
Equations \cref{eq:force,eq:mag} then give identical force and magnetic
measurements; a constant $Q$ also preserves body angular rate. At constant
$R$ and $\vc a=0$, choosing
$\vc d=R\tr(Q\tr\vc g-\vc g)$ gives another stationary history with
constant bias. This is a genuine local quiet sensor ambiguity, despite a
three-dimensional field.

It does \emph{not} imply that fixed heading makes horizontal bias
unobservable. With roll/pitch excitation and a constant bias, the same
$\vc d$ generally cannot satisfy \cref{eq:gauge} at every pose. Lie
derivatives illustrate the mechanism:
\begin{align}
 L_fh_m&=-\sk{\vc\omega_m-\vc b_g}R\tr\vc B,\\
 L_fh_a&=-\sk{\vc\omega_m-\vc b_g}R\tr(\vc a-\vc g)
          +R\tr\dot{\vc a}+\dot{\vc b}_a.
\end{align}
The stationary constant-bias local observability matrix for
$(\phi,b_g,b_a)$ has rank eight of nine when gravity and field are
nonparallel. Its null direction is explicitly checked. Multiple tilting
poses can remove that direction without yaw changes.

If slow physical bias histories are allowed, however,
$\vc d(t)=R(t)\tr(Q\tr\vc g-\vc g)$ gives $\vc a'=Q\vc a$.
Bounded wave displacement remains bounded after rotation. This is a
sensor-level compatibility family, not a theorem of stochastic
non-identifiability: an OU bias law, CARMA excitation anisotropy, calibrated
reference prior or qualified hull-response relationship may assign the
histories different probabilities. It explains why process-prior assumptions
must remain visible in any practical observability claim.

\subsection{Nuisance-eliminated finite-horizon information}
The independent diagnostic builds a scaled sensor/process Jacobian over a
60-second window. Variables include every local pose, all 18 motion states
at every node, constant gyro and accelerometer biases, and optionally a
magnetic-reference perturbation. CARMA transition/noise matrices come from
the existing filter export. The window uses one-second nodes; endpoint
quadrature of nominal attitude over a node is an explicit approximation,
separate from the live 200 Hz integration.

For target columns $A$ and nuisance columns $B$, information is
\begin{equation}
 \mathcal I=A\tr(I-BB^\dagger)A.
\end{equation}
Sparse least-squares projection eliminates nuisance states without a
physical ridge prior. Numerical convergence, scaling and covariance floors
are recorded. Three magnetic assumptions are separated: fixed reference,
an uncertain constant reference with a stated Gaussian prior, and a free
reference. Sensor-only rows deliberately free the motion trajectory; their
bias-information loss is a diagnostic relaxation. Adding CARMA process rows
and the stationary motion prior quantifies model-supplied information.
Varying that prior probes practical weakness. No window result feeds the
live state or covariance.

\subsection{Shared directional gain veto}
An inexpensive candidate monitor accumulates a fading-memory conditional
six-dimensional information matrix for constant world attitude error and
body bias. Its force sensitivity is
\begin{equation}
 \begin{aligned}
 H_a&=[\sk{\vc g}s_\phi\quad -Rs_b]/\sigma_{a,\mathrm{eff}},\\
 H_m&=[-\sk{\vc B}s_\phi\quad0]/\sigma_m.
 \end{aligned}
\end{equation}
The conditional model assumes a known reference and approximately
zero-mean wave acceleration. Motion power inflates force variance; an
explicit effective correlation time prevents 200 Hz copies being counted
as independent stationary poses. Defaults are a 30 s fading horizon, 0.1 s
commit, 1 s correlation time, scales $s_\phi=0.05$ rad and
$s_b=0.05$ m/s$^2$. The force floor is 0.15 m/s$^2$ and the field scale
8 $\mu$T. These values are fixed before holdout evaluation.

After eliminating the three attitude nuisance coordinates, let
$\mathcal I_b=U\operatorname{diag}(\lambda_i)U\tr$. The body-bias gain
multiplier is
\begin{equation}
 D=U\operatorname{diag}\!\left(\frac{\lambda_i}{\lambda_i+1}\right)U\tr,
 \qquad 0\preceq D\preceq I.
\end{equation}
QMEKF uses $K_b\leftarrow DK_b$; TFG/TAN use the equivalent world-bias
factor $RDR\tr$. Other gain rows are retained. Native holding and process
noise remain authoritative. This monitor is \emph{not} the full marine
observability Gramian; it neglects nuisance evolution and cannot certify a
learnable direction. The independent window analysis tests that limitation.

The complete Joseph formula is essential for a modified gain $L$:
\begin{align}
 P_L&=(I-LH)P(I-LH)\tr+L\mathcal R L\tr,\\
 P_L-P_K&=(L-K)(HPH\tr+\mathcal R)(L-K)\tr\succeq0.
\end{align}
Thus, under the stated linear likelihood, gain suppression never gains
information relative to the ordinary Kalman update. This identity is
tested. The monitor contributes no pseudo-measurement. It does not by itself
prove consistency under colored noise, reference error, prior mismatch or
nonlinear corrections. A lower bias RMS with dishonest uncertainty is not
an adoption result.

\section{Controlled Experiments}
The new engine reuses the repository's independent directional spectral
vessel-response truth and sensor-noise generators. It never generates
truth from a filter estimate or feeds truth into normal estimation.
Every compared variant receives one identical timestamped sensor history
per seed. Hashes bind the actual input stream. The gyro is an instantaneous
sample at the force timestamp; current main's endpoint integrator is used.
For variable intervals the elapsed previous interval is passed to prediction.
Prediction, accelerometer correction and due magnetic packets follow the
existing fusion schedule.

Seven primary variants are QMEKF shared modes (QSM), TFG shared modes (TFG),
tangent-bias shared modes (TAN), the same three with the gain veto (-G), and
independent-axis CARMA(4,1). Three additional wave arms enable the veto only
at 300 s (-LG), without resetting adaptation or covariance. The native
front end acquires its reference ordinarily. Exact/wrong reference arms are
explicit diagnostic controls; truth is never a normal update channel.

Development seeds 7301--7302 and short runs precede the fixed holdout seeds
4101--4106. Cases cover quiet water, independent rocking at fixed heading,
low/high waves, rise/recovery, deliberate yaw, horizontal/near-vertical
field, reference error, early and late state perturbations, gyro error,
engine vibration, colored noise, slow accelerometer/gyro bias drift, two-hour motion error, 100/200/400 Hz IMU,
10/25/100 Hz magnetometer and nonuniform sampling. Early perturbations
occur at the first native Live sample; they test convergence after the
shared startup, not a different alignment algorithm. Late perturbations
occur at 300 s. Additional wave arms match synthetic white-noise models,
then separately the population bias prior. Tuning is not optimized against
these holdout histories.

The default directional JONSWAP sea has $H_s=0.8$ m,
$f_p=0.29$ Hz and peak enhancement 3.3. Low/high controls use
$H_s=0.15/4$ m. Transition envelopes multiply wave amplitude
$1\to4\to1$, with 30 s smooth transitions centered at 300/600 s;
their derivatives enter the independent physical truth. Default magnetic
strength/inclination are 50 $\mu$T/65 degrees. White packet standard
deviations are 0.0148 m/s$^2$, 0.00157 rad/s and 0.8 $\mu$T.
Independent per-axis turn-on biases are uniform within
$\pm0.008g$ and $\pm0.05$ degrees/s. The population control matches
the variance of that uniform distribution, not its entire distribution.
Engine disturbances add 0.08/0.15 m/s$^2$ sinusoids at 37/43 Hz
to body X/Z; the colored force control has 0.025 m/s$^2$ stationary
standard deviation and 0.3 s correlation time.

The additional drift controls use the existing physical-history generator's
sinusoidal/ramp bias channels, not estimated states. Rate controls retain
per-packet white-noise standard deviations; they do not model an identical
analog noise spectrum after different anti-alias filters.

Full-state predictions/corrections run at sensor rate. Scores use 20 Hz
estimates, with common windows. Traces and covariance audits are retained at
1 Hz. Metrics include attitude RMS/maximum, velocity by axis, heave,
horizontal wave-position error change, heave-error slope, all six bias
components, first sustained sampled convergence, innovations, numerical failures,
CPU update cost and object storage. Plots use the established shared vector
export/decimation path and are never scoring input.

NEES uses common errors
$[\log(\hat RR\tr),\hat p-p,\hat v-v,\hat b_g-b_g,\hat b_a-b_a]$ and the
full transported 15-dimensional marginal covariance. Window NIS means use
actual accepted measurements, with accepted fractions retained in the CSV;
fixed-time packet NIS is taken before the accept/reject decision. At a fixed timestamp, independent-seed
ensemble bounds use $\chi^2_{nd}$ for $n$ seeds and $d$ dimensions. Serial
samples do not multiply the nominal degrees of freedom. Paired Student
intervals use seed-level changes, not differences of marginal intervals.
Innovation autocorrelation/whiteness tests are diagnostics of the stated
likelihood. Native model mismatch is kept visible.
The chi-square bounds are nominal Gaussian reference bounds, not
distribution-free coverage guarantees under nonlinear dynamics or uniform
turn-on biases. Paired Student intervals also have limited small-sample
precision with six seeds. Conclusions do not rely on treating correlated
time samples or repeated conditions as independent trials.

\section{Measured Results and Ablations}
\StudyCohort
\PrimaryResultsTable
\BiasResultsTable
\PairedResultsTable
\ConsistencyResultsTable
\RecoveryResultsTable
\RobustnessResultsTable
\CostResultsTable
\TimingResultsTable
\AblationResultsTable
\StudyInterpretation
Excess state NEES and small packet NIS are separate failures of the stated
uncertainty model. A broad measurement likelihood can coexist with missing
state uncertainty from magnetic reference and low-frequency motion. Small
NIS is therefore not scored as an improvement. The gate's weaker state
overconfidence still lies outside the fixed-time ensemble bound.
Zero numerical failures also does not mean robust accuracy. High-sea and
near-vertical-field trials show large attitude excursions. In two of the
six 87-degree-field histories, native bias learning remains held throughout
the scored window, with roughly one third of magnetic packets rejected.
The prototype preserves that hold and uses TFG geometry in that interval;
it therefore cannot claim a tangent-bias recovery in those histories.
This exposes a startup/service limitation under weak magnetic geometry,
not a general claim that constant-heading bias estimation is impossible.
The separate pinned vessel-response replay uses the same eight records,
sampling contract and unchanged TFG numeric gates. The documentation build
regenerates its single-realization table from the current executable;
these rows do not replace the paired Monte Carlo study.
\input{tangent-replay-generated.tex-part}

\begin{figure*}[t]\centering
\IncludePGF{tangent_accuracy.pgf}
\caption{Holdout comparisons by physical condition. Curves summarize measured
seed-level errors; uncertainty is across independent seeds. Lower error
alone does not establish an acceptable estimator.}\end{figure*}
\begin{figure*}[t]\centering
\IncludePGF{tangent_bias_tracking.pgf}
\caption{Representative accelerometer-bias drift (left) and gyroscope-bias drift
(right). Each column uses its own declared physical history; all methods
receive identical sensor streams within that history.}\end{figure*}
\begin{figure*}[t]\centering
\IncludePGF{tangent_motion.pgf}
\caption{One-Hz visualization of attitude error and wave-relative motion,
with a final-seconds detail. Faster spectral content is not resolved by
these display traces; RMS scoring uses 20 Hz estimates from full-rate updates.
Horizontal displacement is not absolute geographic position.}\end{figure*}
\begin{figure*}[t]\centering
\IncludePGF{tangent_degradation.pgf}
\caption{Engine-noise and sea-state rise/recovery controls. Common noise and
truth isolate estimator behavior under each disturbance.}\end{figure*}
\begin{figure*}[t]\centering
\IncludePGF{tangent_consistency.pgf}
\caption{Consistency and innovation diagnostics. Fixed-time ensemble bounds
use independent seeds. Serial coverage is not a calibration certificate.}\end{figure*}
\begin{figure*}[t]\centering
\IncludePGF{tangent_observability.pgf}
\caption{Nuisance-eliminated local information under explicit motion and
magnetic-reference assumptions. Numerical rank thresholds are diagnostic,
not a global observability proof.}\end{figure*}
\clearpage\twocolumn

\section{Implementation, Regression and Embedded Feasibility}
The new class inherits the current shared TFG estimator. The existing
minimal motion realization, confidence, period/amplitude adaptation,
magnetic acquisition and timing contract are reused. The $\mathrm{SE}(3)$
Jacobian cross block is extracted once and retained under its former
namespace alias. A shared gain interface supports the optional bias veto in
both QMEKF and TFG; defaults disable it. Public covariance storage and
ordinary estimator APIs remain compatible.

The tangent correction uses fixed-size small matrices and rank-row
congruences. The covariance remains 27 by 27 (2916 bytes in float) and no
additional full covariance scratch is introduced. The optional monitor
stores a 6 by 6 information matrix plus small matrices. Wrapper storage
includes its heap-owned estimator; reported sizes exclude allocator
metadata and are host ABI values.
The common optional gain/monitor interface adds 320 host bytes to each
shared-mode wrapper/core even when disabled, relative to the untouched-main
build of the same harness. The independent-axis core's interface adds
32 host bytes. These costs are retained in the baseline comparison.
CPU timings compare the same host build
and instrument only filter updates. They are not Xtensa execution or board
latency measurements; stack high-water marks, firmware linking and 200 Hz
headroom on ESP32-S3 remain separate hardware qualifications.

Regression checks cover group laws, the adjoint, exponential/logarithm,
finite-difference reset, structured versus dense covariance transport,
positive semidefiniteness, quiet/rocking information, Joseph gain penalties,
and complete baseline traces built against an untouched main worktree.
The original TFG replay limits are shared as one fixture; experimental
replay failures are reported rather than used to widen them. The root
\texttt{make all} and existing unrelated replay suites retain their original
contracts. A focused CI job builds the implementation and article and runs
independent smoke histories. The full cohort is reproducible separately.
Separate checks confirm identical physical estimates with/without the gate
while the native bias hold is active, and before late gate activation at
300 s. Consequently the late-activation comparison measures live learning
without changing startup estimates.
\ValidationResults

\section{Recommendation and Limitations}
\StudyRecommendation
The scope remains a local stochastic marine estimator. The conditional
gain monitor is not a deployable observability certificate; a reliable
sensor-driven gate must account for magnetic calibration/reference
uncertainty and correlated low-frequency motion. A more accurate prior,
independently qualified attitude, confirmed stationary multi-pose data,
time-aligned GNSS/DVL velocity or pressure would be additional information,
and must be evaluated as such rather than credited to a group choice.

A full 30-state TG-INS with a compatible stochastic augmentation, an
iterated measurement objective and realistic recorded hardware data are
future experiments. This study neither implements that full-TG comparator
nor establishes that all tangent-group filters are inferior. It tests the
27-state model-preserving tangent-bias alternative and separates it from a
shared learning policy. Statistical acceptance requires a paired accuracy
gain, adequate uncertainty, robustness and acceptable target cost together.
Finite Monte Carlo runs do not prove nonlinear global stability.

\begin{thebibliography}{9}
\bibitem{fornasier2022} A. Fornasier, Y. Ng, R. Mahony and S. Weiss,
``Equivariant Filter Design for Inertial Navigation Systems with Input
Measurement Biases,'' ICRA, 2022, doi:10.1109/ICRA46639.2022.9811778.
\bibitem{barrau2023} A. Barrau and S. Bonnabel, ``The Geometry of Navigation
Problems,'' IEEE Trans. Automatic Control, vol. 68, no. 2, pp. 689--704,
2023, doi:10.1109/TAC.2022.3144328.
\bibitem{fornasier2025} A. Fornasier, Y. Ge, P. van Goor, R. Mahony and
S. Weiss, ``Equivariant symmetries for inertial navigation systems,''
Automatica, vol. 181, 112495, 2025, doi:10.1016/j.automatica.2025.112495.
\bibitem{fornasier2024} A. Fornasier, ``Equivariant Symmetries for Aided
Inertial Navigation,'' dissertation, 2024, arXiv:2407.14297.
\bibitem{vantoor2026} P. van Goor and J. R. Forbes, ``The Iterative
Equivariant Filter,'' preprint, September 2026, arXiv:2609.12328.
\bibitem{lunawat2026} A. Lunawat, P. van Goor, F. Dellaert and S. B.
Williams, ``Equivariant Filter Design for Acoustic and Depth Aided Inertial
Navigation Systems,'' preprint, September 2026, arXiv:2609.19742.
\end{thebibliography}
\end{document}
